% !TeX root = main.tex
% !TeX spellcheck = pt_BR

\chapter{Inferência Bayesiana}

\begin{theorem}[Normal com $\sigma^2$ conhecido]
	\leavevmode
	Seja $\v{x} = (x_1, \ldots, x_n)$.
	\begin{outline}[enumerate]
		\1 Verosimilhança: \\
		\[ (X_i|\mu) \sim \normdist{\mu,\sigma^2} \] 
		
		\1 Priori:
		\[ \mu \sim \normdist{\mu_0, \sigma_0^2} \]
		
		\1 Posteriori:
		\[
		(\mu | \v{x}) \sim \normdist{\mu_1, \sigma_1^2}
		\]
		\begin{align}
			\sigma_1^2 &= 	\frac{1}{\frac{n}{\sigma^2}+\frac{1}{\sigma_0^2}}= \frac{\sigma^2\sigma_0^2}{n\sigma_0^2+\sigma^2} \quad \\
			\mu_1 &= \sigma_1^2\left( \frac{n\bar{x}}{\sigma^2} + \frac{\mu_0}{\sigma_0^2}  \right) = 	\frac{\frac{n\bar{x}}{\sigma^2}+\frac{\mu_0}{\sigma_0^2}}{\frac{n}{\sigma^2}+\frac{1}{\sigma_0^2}}= \frac{n\bar{x}\sigma_0^2+\mu_0\sigma^2}{n\sigma_0^2+\sigma^2}
		\end{align}
		
		Fazendo $\tau=1/\sigma^2$:
		\begin{align}
			\tau_1 &= \tau_0 + n\tau \\
			\mu_1 &= \frac{n\tau\bar{x}+\tau_0\mu_0}{n\tau+\tau_0}
		\end{align}
		
		\begin{note}
			A média a posteriori é uma combinação convexa da média a priori e da média amostral.
		\end{note}
	\end{outline}
\end{theorem}

\begin{theorem}[Normal multivariada com $\m{\Sigma}$ conhecido]
	\leavevmode
	Seja $\v{y} = (y_1, \ldots, y_n)$.
	\begin{outline}[enumerate]
		\1 Verosimilhança:
		\[ p(\v{y}|\v{\mu}) = \prod_{n=1}^{N} \normdist{y_n \mid \v{\mu}, \m{\Sigma}} \sim \normdist{\bar{\v{y}} \mid \v{\mu}, \frac{1}{N}\m{\Sigma}} \] 
		
		\1 Priori:
		\[ p(\v{\mu}) \sim \normdist{\v{\mu} \mid \prior{\v{m}}, \prior{\m{V}}} \]
		
		\1 Posteriori:
		\[
		p(\v{\mu} | \v{y}) \sim \normdist{\v{\mu} \mid \post{\v{m}}, \post{\m{V}}}
		\]
		\begin{align}
			\post{\m{V}}^{-1} &= \prior{\m{V}}^{-1} + N \m{\Sigma}^{-1} \\
			\post{\v{m}} &= \post{\m{V}} \left[ \m{\Sigma}^{-1} (N\bar{\v{y}}) + \prior{\m{V}}^{-1} \prior{\v{m}} \right]  
		\end{align}
	\end{outline}
\end{theorem}

\begin{theorem}[Regressão linear bayesiana com $\m{\Sigma}$ conhecido]
	\leavevmode
	Seja $\v{y} = (y_1, \ldots, y_n)$.
	\begin{outline}[enumerate]
		\1 Verosimilhança:
		\[ p(\v{y}|\v{\mu}) \sim \normdist{\bar{\v{y}} \mid \m{\Fcal}\v{\mu}, \m{\Sigma}} \] 
		
		\1 Priori:
		\[ p(\v{\mu}) \sim \normdist{\v{\mu} \mid \prior{\v{m}}, \prior{\m{V}}} \]
		
		\1 Posteriori:
		\[
		p(\v{\mu} | \v{y}) \sim \normdist{\v{\mu} \mid \post{\v{m}}, \post{\m{V}}}
		\]
		\begin{align}
			\post{\m{V}}^{-1} &= \prior{\m{V}}^{-1} + \m{\Fcal}\T \m{\Sigma}^{-1} \m{\Fcal} \\
			\post{\v{m}} &= \post{\m{V}} \left[ \m{\Fcal}\T  \m{\Sigma}^{-1}{\v{y}} + \prior{\m{V}}^{-1} \prior{\v{m}} \right]  
		\end{align}
	\end{outline}
\end{theorem}

\begin{theorem}[Distribuição condicional e regressão linear]
	Suponha que o vetor-$p$ $\v{Y}$ e o vetor-$n$ $\v{\theta}$ estejam relacionados via a distribuição condicional
	\[
	(\v{Y} \mid \v{\theta}) \sim \normdist{\m{F}'\v{\theta}, \m{V}},
	\]
	onde a matriz $(n \times p)$ $\m{F}$ e a matriz $(p \times p)$ definida positiva simétrica $\m{V}$ são constantes, e que a distribuição marginal de $\v{\theta}$ seja dada por
	\[
	\v{\theta} \sim \normdist{\v{a}, \m{R}}.
	\]
	
	Segue que
	\[
	\begin{bmatrix}
		\v{Y} \\ \v{\theta}
	\end{bmatrix} \sim \normdist{\begin{bmatrix}
			\m{F}'\v{a} \\ \v{a}
		\end{bmatrix}, \begin{bmatrix}
			\m{F}'\m{R}\m{F} + \m{V} & \m{F}'\m{R} \\
			\m{R}\m{F} & \m{R}
	\end{bmatrix}}.
	\]
	
	Então $(\v{\theta} \mid \v{Y}) \sim \normdist{\v{m}, \m{C}}$, onde
	\[
	\v{m} = \v{a} + \m{R}\m{F}[\m{F}'\m{R}\m{F} + \m{V}]^{-1}[\v{Y} - \m{F}'\v{a}]
	\]
	e
	\[
	\m{C} = \m{R} - \m{R}\m{F}[\m{F}'\m{R}\m{F} + \m{V}]^{-1}\m{F}'\m{R}.
	\]
	
	Definindo $\v{e} = \v{Y} - \m{F}'\v{a}$, $\m{Q} = \m{F}'\m{R}\m{F} + \m{V}$ e $\m{A} = \m{R}\m{F}\m{Q}\inv$, estas equações tornam-se
	\[
	\v{m} = \v{a} + \m{A}\v{e} \quad \text{e} \quad \m{C} = \m{R} - \m{A}\m{Q}\m{A}'.
	\]
\end{theorem}

\begin{theorem}[Normal-Gamma]
	\leavevmode 
	\begin{outline}
		\1 Verosimilhança: 
		\[ Y_1, \ldots, Y_k \sim \normdist{\mu, \phi^{-1}}\]
		
		\1 Priori:
		\begin{align}
			(\mu,\phi) &\sim \ngdist{m, C, n, d} \\
			\mu|\phi &\sim \normdist{m, C\phi^{-1}} \\
			\phi &\sim \gammadist{n/2, /2}
		\end{align}
		
		\1 Posteriori: \[ \ (\mu,\phi)|\v{y} \sim \ngdist{m^*, C^*, n^*, d^*}\]
		\begin{align}
			n^* &= n+k \\
			m^* &= \frac{km/C+\sum_{i=1}^{k}y_i}{k+1/C} \\
			C^* &= \frac{C}{2+kC} \\
			d^* &= d + \sum_{i=1}^{k}(y_i-\bar{y})^2 + \frac{k(\bar{y}-m)^2}{1+kC}
		\end{align}	
		
		\1 Marginais:
		\begin{align}
			\mu &\sim \tdist{n}(m,R), \quad R=C\frac{d}{n} \\
			\mu|\v{y} &\sim \tdist{n}(m^*, R^*), \quad R^*=C^*\frac{d^*}{n^*}
		\end{align}
	\end{outline}
\end{theorem}

\begin{theorem}[Normal-Gamma Multivariada]
	\leavevmode 
	\begin{outline}
		\1 Verossimilhança: 
		\[ \v{Y}_1, \ldots, \v{Y}_k \mid \v{\theta}, \phi \stackrel{iid}{\sim} \normdist{\v{\theta}, \phi^{-1}\v{I}_p}\]
		
		\1 Priori:
		\begin{align}
			(\v{\theta},\phi) &\sim \ngdist{\v{m}, \v{C}, n, d} \\
			\v{\theta}|\phi &\sim \normdist{\v{m}, \v{C}\phi^{-1}} \\
			\phi &\sim \gammadist{n/2, d/2}
		\end{align}
		
		\1 Posteriori: 
		\[ (\v{\theta},\phi)|\v{Y}_1, \ldots, \v{Y}_k \sim \ngdist{\v{m}^*, \v{C}^*, n^*, d^*}\]
		\begin{align}
			n^* &= n+kp \\
			\v{m}^* &= (\v{C}^{-1} + k\v{I}_p)^{-1}(\v{C}^{-1}\v{m} + k\bar{\v{Y}}) \\
			\v{C}^* &= (\v{C}^{-1} + k\v{I}_p)^{-1} \\
			d^* &= d + \sum_{j=1}^k (\v{Y}_j - \bar{\v{Y}})'(\v{Y}_j - \bar{\v{Y}}) + (\bar{\v{Y}} - \v{m})'(\v{C}^{-1} + k\v{I}_p)^{-1}(\bar{\v{Y}} - \v{m})
		\end{align}	
		
		\2 Caso particular $k=1$: Após observar $\v{Y}_1 = \v{Y}$,
		\begin{align}
			n^* &= n+p \\
			\v{m}^* &= (\v{C}^{-1} + \v{I}_p)^{-1}(\m{C}^{-1}\v{m} + \v{Y}) \\
			\m{C}^* &= (\m{C}^{-1} + \v{I}_p)^{-1} \\
			d^* &= d + (\v{Y} - \v{m})'(\m{C}^{-1} + \v{I}_p)^{-1}(\v{Y} - \v{m})
		\end{align}
		
		\1 Marginais:
		\begin{align}
			\v{\theta} &\sim \text{T}_n[\v{m}, \v{R}], \quad \v{R}=\m{C}\frac{d}{n} \\
			\v{\theta}|\v{Y}_1, \ldots, \v{Y}_k &\sim \text{T}_{n^*}[\v{m}^*, \v{R}^*], \quad \v{R}^*=\v{C}^*\frac{d^*}{n^*}
		\end{align}
		
		\1 Marginais univariadas: Se $\theta_i$ é o $i$-ésimo elemento de $\v{\theta}$, então
		\begin{align}
			\theta_i &\sim \tdist{n}(m_i, R_{ii}), \quad R_{ii}=C_{ii}\frac{d}{n} \\
			\theta_i|\v{Y}_1, \ldots, \v{Y}_k &\sim \tdist{n^*}(m_i^*, R_{ii}^*), \quad R_{ii}^*=C_{ii}^*\frac{d^*}{n^*}
		\end{align}
		onde $\bar{\v{Y}} = \frac{1}{k}\sum_{j=1}^k \v{Y}_j$
	\end{outline}
	
	\begin{note}
		\leavevmode
		\begin{outline}
			\1 O vetor de parâmetros $\v{\theta}$ tem dimensão $p \times 1$.
			\1 Cada observação $\v{Y}_j$ também tem dimensão $p \times 1$.
			\1 $\m{C}$ é uma matriz $p \times p$ simétrica positiva definida.
			\1 Os graus de liberdade aumentam de $n$ para $n^* = n + kp$ (cada observação vetorial contribui com $p$ graus de liberdade).
			\1 A notação $\text{T}_n[\v{m}, \m{R}]$ denota a distribuição t multivariada com $n$ graus de liberdade, moda $\v{m}$ e matriz de escala $\m{R}$.
		\end{outline}
	\end{note}
\end{theorem}

\begin{theorem}[Normal-Gamma para Regressão Linear]
	\leavevmode 
	\begin{outline}
		\1 Verossimilhança: 
		\[ (Y_i \mid \v{\theta}, \phi) \sim N[\v{F}_i'\v{\theta}, k_i\phi^{-1}], \quad i = 1, \ldots, n\]
		
		\1 Priori:
		\begin{align}
			(\v{\theta},\phi) &\sim \text{NG}(\v{a}, \v{R}, \nu, s) \\
			\v{\theta}|\phi &\sim N[\v{a}, \v{R}\phi^{-1}] \\
			\phi &\sim \text{Gamma}(\nu/2, s/2)
		\end{align}
		
		\1 Posteriori para múltiplas observações: 
		\[ (\v{\theta},\phi)|Y_1, \ldots, Y_n \sim \text{NG}(\v{a}^*, \v{R}^*, \nu^*, s^*)\]
		\begin{align}
			\nu^* &= \nu + n \\
			Q_i &= \v{F}_i'\v{R}\v{F}_i + k_i \quad \text{(variância preditiva)} \\
			e_i &= Y_i - \v{F}_i'\v{a} \quad \text{(erro de predição)} \\
			\v{R}^* &= \v{R} - \sum_{i=1}^n \frac{\v{R}\v{F}_i\v{F}_i'\v{R}}{Q_i} \\
			\v{a}^* &= \v{a} + \sum_{i=1}^n \frac{\v{R}\v{F}_i e_i}{Q_i} \\
			s^* &= s + \sum_{i=1}^n \frac{e_i^2}{Q_i}
		\end{align}	
		
		\2 Caso particular $n=1$: Após observar $(Y_1, \v{F}_1)$ com $Y_1 = Y$, $\v{F}_1 = \v{F}$, $k_1 = k$,
		\begin{align}
			\nu^* &= \nu + 1 \\
			f &= \v{F}'\v{a} \quad \text{(predição priori)} \\
			Q &= \v{F}'\v{R}\v{F} + k \quad \text{(variância preditiva priori)} \\
			\v{R}^* &= \v{R} - \frac{\v{R}\v{F}\v{F}'\v{R}}{Q} \\
			\v{a}^* &= \v{a} + \frac{\v{R}\v{F}(Y - f)}{Q} \\
			s^* &= s + \frac{(Y - f)^2}{Q}
		\end{align}
		
		\3 Distribuições condicionais e marginal posteriori ($n=1$):
		\begin{align}
			(\v{\theta} \mid \phi, Y) &\sim N[\v{a}^*, \v{R}^*\phi^{-1}] \\
			(\phi \mid Y) &\sim \text{Gamma}\left(\frac{\nu+1}{2}, \frac{s + (Y-f)^2/Q}{2}\right)
		\end{align}
		
		\1 Marginais:
		\begin{align}
			\v{\theta} &\sim \text{T}_\nu[\v{a}, \v{R}S], \quad S = s/\nu \\
			\v{\theta}|Y_1, \ldots, Y_n &\sim \text{T}_{\nu^*}[\v{a}^*, \v{R}^*S^*], \quad S^* = s^*/\nu^*
		\end{align}
		
		\1 Marginais univariadas: Se $\theta_j$ é o $j$-ésimo elemento de $\v{\theta}$, então
		\begin{align}
			\theta_j &\sim t_\nu\left(a_j, R_{jj}S\right) \\
			\theta_j|Y_1, \ldots, Y_n &\sim t_{\nu^*}\left(a_j^*, R_{jj}^*S^*\right)
		\end{align}
		ou de forma padronizada:
		\begin{align}
			\frac{\theta_j - a_j}{\sqrt{R_{jj}S}} &\sim t_\nu(0, 1) \\
			\frac{\theta_j - a_j^*}{\sqrt{R_{jj}^*S^*}} &\sim t_{\nu^*}(0, 1)
		\end{align}
		
		\1 Distribuição preditiva: Para uma nova observação $Y_{n+1}$ com covariáveis $\v{F}_{n+1}$ e fator $k_{n+1}$,
		\begin{align}
			Y_{n+1}|Y_1, \ldots, Y_n &\sim t_{\nu^*}\left(f^*, Q^*S^*\right) \\
			f^* &= \v{F}_{n+1}'\v{a}^* \\
			Q^* &= \v{F}_{n+1}'\v{R}^*\v{F}_{n+1} + k_{n+1}
		\end{align}
	\end{outline}
	
	\begin{note}
		\leavevmode
		\begin{outline}
			\1 O vetor de parâmetros $\v{\theta}$ tem dimensão $p \times 1$.
			\1 Cada $Y_i$ é uma observação escalar com vetor de covariáveis $\v{F}_i$ ($p \times 1$).
			\1 $k_i$ é o fator de escala da variância da $i$-ésima observação.
			\1 $\v{R}$ é uma matriz $p \times p$ simétrica positiva definida.
			\1 $S = s/\nu = 1/E[\phi]$ é o inverso da precisão esperada.
			\1 Os graus de liberdade aumentam de $\nu$ para $\nu^* = \nu + n$ (número de observações escalares).
			\1 A notação $\text{T}_\nu[\v{a}, \v{R}S]$ denota a distribuição t multivariada com $\nu$ graus de liberdade, vetor de localização $\v{a}$ e matriz de escala $\v{R}S$.
			\1 A atualização sequencial é possível: a posteriori após $i$ observações serve como priori para a $(i+1)$-ésima observação.
		\end{outline}
	\end{note}
\end{theorem}